\chapter*{A Map of Protocols and Acronyms}
\addcontentsline{toc}{chapter}{A Map of Protocols and Acronyms}
The same six-letter families recur throughout this book; pin them down once
here so later chapters can be read without re-decoding:

\begin{center}
\begin{tabular}{lp{4.2in}}
\textbf{Term} & \textbf{Meaning in this book} \\
\hline
TM & Transactional Memory (the concept, Ch.~1--2) \\
STM & Software TM: conflict detection in the runtime (Part~IV) \\
HTM & Hardware TM: conflict detection in the cache hierarchy (Ch.~16) \\
SGL & Single Global Lock --- the reference non-speculative design (Ch.~8) \\
TL2 & Global-clock, versioned-lock STM (Ch.~9) \\
NOrec & No Read-Lock, value-based validation (Ch.~9) \\
TinySTM & Lock-based STM with per-object locks (WBCTL/WT variants, Ch.~9) \\
JVSTM / MVLog & Multi-version STMs (Ch.~9) \\
PlusCal / TLA+ / TLC & Specification language, logic, and model checker (Ch.~6) \\
OCC / read-validate & Optimistic concurrency control; validate-then-write-back (Ch.~9) \\
RTM & Intel Restricted Transactional Memory (Ch.~16) \\
GUST & AtomicINC commit-log GPU TM (Ch.~17) \\
\end{tabular}
\end{center}

\section*{Why the book is ordered this way}
\begin{itemize}
    \item The order is deliberately \emph{not} ``hardware then software then
          formal then compilers.'' Mathematics and verification (Part~III)
          precede the runtimes (Part~IV) they specify, so that ``what a
          correct runtime must preserve'' is understood before ``what the
          runtime does with the calls.'' The runtimes in turn precede the
          compiler instrumentation (Part~V) that targets them, so the runtime
          ABI of Chapter~7 is fixed before ``what an instrumentation pass
          must preserve'' is read.
    \item Readers who prefer runtimes-first can start at Part~IV and consult
          Parts~III and V on demand; cross-references mark every such jump.
\end{itemize}

